# Methods Supplementary Table S2 — Rare Immune Subpopulation Catalog

**Source:** `workspace/registry_v1/populations.tsv` (stable IDs) + `workspace/marker_panels.md` (full panels).
**Format target:** LaTeX longtable, drops into manuscript supplementary as one floating table per page.
**Status:** ready for Biological Science's bib compile.

```latex
\begin{longtable}{@{}l l l p{32mm} p{30mm} c l@{}}
\caption{Rare immune subpopulations evaluated in this study. Prevalence is reported as fraction of the relevant parent compartment. IF\textsubscript{score} is the integration-fragility score (see Methods \ref{sec:if-score}). Fragility mode follows the revised three-object taxonomy of Methods \ref{sec:loss-mode-taxonomy}: A, local-density log-ratio; B, cross-batch persistent-homology bottleneck; C, spectral Laplacian gap. Clinical weight is an ordinal on $\{0.4,0.7,1.0\}$ (mechanistic, biomarker, therapy target). Full marker panels, including supportive and exclusion gates, are in Supplementary Table S3.} \label{tab:s2-rare-populations}\\
\toprule
Pop ID & Name & Compartment & Core markers (positive) & Adjacent abundant neighbor & Prevalence & IF\textsubscript{score} \\
\midrule
\endfirsthead
\multicolumn{7}{c}{\tablename\ \thetable\ continued from previous page} \\
\toprule
Pop ID & Name & Compartment & Core markers (positive) & Adjacent abundant neighbor & Prevalence & IF\textsubscript{score} \\
\midrule
\endhead
\midrule
\multicolumn{7}{r}{continued on next page} \\
\endfoot
\bottomrule
\endlastfoot
IM-T-001 & TPEX (CXCR5+TCF7+)        & CD8 T         & \textit{TCF7, CXCR5, IL7R}                & Tex\textsubscript{term}                & 0.05--0.3\%  & 0.88 \\
IM-T-002 & TRM-TR (CD103+CD39+)      & CD8 T         & \textit{CD69, ITGAE, ENTPD1}              & Tex\textsubscript{effector}            & 0.3--1\%     & 0.82 \\
IM-T-003 & eTreg (CCR8+TNFRSF9+)     & Treg          & \textit{FOXP3, IL2RA, CTLA4, CCR8}        & Resting Treg                           & 0.1--0.5\%   & 0.90 \\
IM-T-004 & T\textsubscript{FH}/T\textsubscript{FR} (TLS) & CD4 T & \textit{CXCR5, BCL6, PDCD1}         & T\textsubscript{PH} / Tex\textsubscript{CD4} & 0.5--3\%     & 0.72 \\
IM-T-005 & MAIT (tissue)             & CD8\,T unc.   & \textit{KLRB1, SLC4A10, ZBTB16}           & CD8 effector                           & 0.5--5\%     & 0.85 \\
IM-T-006 & $\gamma\delta$-T V$\delta1$ tumor & $\gamma\delta$T & \textit{TRDC, TRGC1, TRDV1}           & NK / CD8 effector                      & 0.05--0.5\%  & 0.78 \\
IM-M-001 & AS DC / DC5 (AXL+SIGLEC6+) & DC           & \textit{AXL, SIGLEC6}                     & pDC / cDC2                             & 0.02--0.1\%  & 0.90 \\
IM-M-002 & pDC                       & DC            & \textit{CLEC4C, LILRA4, IL3RA, TCF4, IRF7}& Naïve B / cDC2                         & 0.05--0.5\%  & 0.80 \\
IM-M-003 & LAMP3+ mregDC             & DC            & \textit{LAMP3, CCR7, BIRC3}               & Generic activated DC                   & 0.02--0.1\%  & 0.85 \\
IM-M-004 & Langerhans (CD207+)       & DC            & \textit{CD207, CD1A}                      & Generic cDC2                           & 0.1--1\%     & 0.70 \\
IM-M-005 & TREM2+ LAM                & Macrophage    & \textit{TREM2, APOE, APOC1}               & FOLR2+ / SPP1+                          & 0.5--3\%     & 0.75 \\
IM-M-006 & FOLR2+ perivascular       & Macrophage    & \textit{FOLR2, LYVE1}                      & Generic TAM                            & 0.05--0.2\%  & 0.78 \\
IM-M-007 & SPP1+ TAM                 & Macrophage    & \textit{SPP1, MARCO}                       & TREM2+ LAM                             & 0.5--3\%     & 0.70 \\
IM-M-008 & ISG-M$\varphi$            & Macrophage    & \textit{IFIT1, IFIT3, ISG15, OAS1}         & SPP1+ TAM                               & 0.2--1\%     & 0.60 \\
IM-B-001 & ABC (T-bet+CD11c+)        & B             & \textit{TBX21, ITGAX, FCRL5}               & Memory B                               & 0.1--5\%     & 0.78 \\
IM-B-002 & TLS memory B / plasma     & B             & \textit{MS4A1, CD27, XBP1, MZB1}           & Generic memory B                       & 1--5\%       & 0.70 \\
IM-N-001 & KIR+ adaptive NK          & NK            & \textit{KLRC2, BCL11B, IFI44L}             & CD56dim canonical NK                   & 0.1--1\%     & 0.75 \\
IM-I-001 & ILC1 (tumor)              & ILC           & \textit{TBX21, IFNG}                       & NK                                     & 0.02--0.2\%  & 0.80 \\
IM-I-002 & ILC2                      & ILC           & \textit{GATA3, PTGDR2, IL1RL1}             & GATA3+ Th                              & 0.01--0.1\%  & 0.75 \\
IM-I-003 & ILC3 (NKp44+)             & ILC           & \textit{RORC, NCR2, AHR}                   & Th17                                   & 0.01--0.1\%  & 0.78 \\
IM-G-001 & Mast cells (tumor)        & Granulocyte   & \textit{TPSAB1, TPSB2, CPA3, KIT, HDC}     & n/a                                    & 0.05--0.5\%  & 0.70 \\
IM-G-002 & LOX-1+ PMN-MDSC           & Granulocyte   & \textit{OLR1, ARG1, CD84}                  & Neutrophil                             & 0.05--0.5\%  & 0.65 \\
IM-H-001 & CD34+ HSPC (contaminating)& Progenitor    & \textit{CD34, KIT, SPINK2}                 & Lymphoid lineage                       & 0.1--0.5\%   & 0.70 \\
EP-E-001 & Pulmonary ionocytes (anchor) & Epithelial & \textit{FOXI1, ASCL3, CFTR}                & Secretory / basal                      & 0.5--2\%     & 0.92 \\
EP-E-002 & Pancreatic epsilon (anchor) & Endocrine   & \textit{GHRL}                              & $\alpha$ / $\beta$                     & 0.05--0.2\%  & 0.75 \\
\end{longtable}
```

## Notes for Biological Science (manuscript lead)

1. This is a LaTeX longtable with 25 rows. If the skeleton uses `tabularx` or `booktabs` without longtable, I can reformat.
2. Full marker panels (positive_core / supportive / exclusion) belong in Supplementary Table S3. I'll draft that next if you confirm format preference (likely the same longtable style).
3. All pop_id values are stable identifiers from `workspace/registry_v1/populations.tsv` — if the registry is published as SN2, this table cross-references it.
4. Prevalence ranges are fraction-of-compartment as specified. For ranges expressed as fraction-of-TME vs fraction-of-parent-lineage, we may want a footnote per row — flag if reviewer will notice.
5. IF_score is the provisional hand-engineered value; will be replaced with Math's principled statistic for the final version.

## Caveats

- Row IM-M-008 (ISG-Mφ) has relatively low IF_score (0.60) because inflammation-state clusters are volatile; preservation success varies strongly with the inflammatory milieu of the input dataset. Flag in caption or move to Tier-2 if the reviewer asks.
- Rows for apCAF (boundary case) and EMT-hybrid (confounder) are omitted from this table — they belong in Methods as methodological caveats, not as rare-immune targets.
- If any row is dropped from the final experiments, also drop it from this table; otherwise the caption commits us to reporting on it.
